\section{Uniform nondegeneracy of the joint return record}
\label{sec:joint-nondegeneracy}
We now combine measurable phase rigidity with the exact integer-valued
section compensation. The proof treats every joint direction of the
same four-coordinate record. It does not pass from the scalar count
variance to a determinant by an inequality.

\subsection{A finite-cover obstruction for displacement}
\begin{lemma}[No real spatial collision coboundary]
\label{lem:no-spatial-coboundary}
If $w\in\R^2$ and $b\in L^2(\nu;\R)$ satisfy
\begin{equation}\label{eq:spatial-coboundary}
 w\cdot\kappa_R=b-b\circ T_R,
\end{equation}
then $w=0$.
\end{lemma}
\begin{proof}
Regard $w$ as a row vector and suppose $w\ne0$.
Composition with $S_1$, using \eqref{eq:collision-rotation-identities},
gives $(wA)\cdot\kappa_R=b\circ S_1-(b\circ S_1)\circ T_R$.
The matrix with rows $w=(w_1,w_2)$ and
$wA=(w_2,-w_1+w_2)$ has determinant
\begin{equation}\label{eq:spatial-rotation-determinant}
 -(w_1^2-w_1w_2+w_2^2),
 \qquad w_1^2-w_1w_2+w_2^2\ge\tfrac12|w|^2>0.
\end{equation}
Solving these two real identities gives a vector-valued
$H\in L^2(\nu;\R^2)$ with $\kappa_R=H-H\circ T_R$.

Consider the physical billiard on the flat torus $\R^2/(2\Lambda)$.
There are four equal scatterers, labeled by
$\ell\in\Lambda/(2\Lambda)\simeq(\Z/2\Z)^2$. In collision
coordinates its map and invariant probability are exactly
\begin{equation}\label{eq:four-sheet-collision}
 \widetilde T_R(x,\ell)
       =(T_Rx,\ell+\kappa_R(x)\bmod2),\qquad
 \widetilde\nu=\nu\otimes\operatorname{unif}((\Z/2\Z)^2).
\end{equation}
This map is ergodic. Indeed the enlarged table has disjoint smooth
strictly convex disks, connected free space and the same finite
horizon; the finite-horizon dispersing-billiard theorem
\cite[Section 8.1, Theorem 6]{Young} applies. To work only with a
rectangular fundamental domain, take the physical cover with periods
$(2,0)$ and $(0,2\sqrt3)$. It has eight scatterers and satisfies the
same hypotheses. Its quotient by the deck group is precisely the
four-scatterer table above, so its ergodicity implies that of
\eqref{eq:four-sheet-collision}. Neither argument changes Euclidean
length or specular reflection.

However,
\[
 F(x,\ell)=\exp\{i\pi(H_1(x)+\ell_1)\}
\]
is invariant under \eqref{eq:four-sheet-collision}: the integer
increment $\kappa_{R,1}$ cancels $H_1\circ T_R-H_1$, and reduction
of $\ell_1$ modulo two does not alter the exponential. The function
has modulus one and integral zero by averaging over $\ell_1$.
It cannot be almost everywhere constant, contradicting ergodicity.
This proves the lemma. No value of $b$ or $H$ at a periodic point is
used.
\end{proof}

\subsection{Removing the section contribution}
\begin{theorem}[Uniform positive definiteness]
\label{thm:joint-uniform-nondegeneracy}
There exist constants $0<d_0\le D_0<\infty$ such that
\begin{equation}\label{eq:joint-uniform-nondegeneracy}
 d_0I_4\le D_R\le D_0I_4\qquad(R\in I).
\end{equation}
Equivalently the actual $L^2$ coboundaries in
Theorem~\ref{thm:gaussian-kernel-coboundary} have only the zero joint
direction. The discontinuity of $\eta_R=\one_{Y_R^*}$ is allowed.
\end{theorem}
\begin{proof}
Fix $R$ and suppose $v^{\mathsf T}D_Rv=0$, where $v=(w,a,b)$ with
$w\in\R^2$. The collision-clock coboundary characterization gives a
real $B\in L^2(\nu)$ such that
\begin{equation}\label{eq:joint-zero-coboundary}
 w\cdot\kappa_R+a+b\tau_R-\zeta\eta_R=B-B\circ T_R,
 \qquad \zeta=\frac{a+b\bar\tau_R}{c_*}.
\end{equation}
All terms are the physical single-collision quantities. We separate
whether the coefficient of the section indicator vanishes.

If $\zeta\ne0$, set $q=\exp(2\pi i B/\zeta)$. Since $\eta_R$ is
integer-valued, exponentiation removes it exactly and yields
\[
 q\circ T_R
 =\exp\!\left(-\frac{2\pi i}{\zeta}
          (w\cdot\kappa_R+a+b\tau_R)\right)q.
\]
Theorem~\ref{thm:physical-phase-rigidity} implies $b=0$ and
$-2\pi a/\zeta\in2\pi\Z$. But now
$\zeta=a/c_*\ne0$, so $a/\zeta=c_*$. Since $0<c_*<1$, this is
impossible. Notice that the rotations are applied only after the
section indicator has disappeared. Their use requires no symmetry of
the actual return section.

If $\zeta=0$, set $q=\exp(iB)$ in
\eqref{eq:joint-zero-coboundary}. Proposition~\ref{prop:measurable-roof-rigidity}
implies $b=0$. The definition of $\zeta$ then gives $a=0$.
The remaining identity is the real spatial coboundary
$w\cdot\kappa_R=B-B\circ T_R$, and
Lemma~\ref{lem:no-spatial-coboundary} gives $w=0$.

Thus $v=0$ in all cases. The symmetric matrix $D_R$ is positive definite
for each $R$. Theorem~\ref{thm:collision-covariance} gives its continuity
on the compact interval $I$. The smallest eigenvalue is consequently
bounded below by a positive number, and the largest is bounded above.
This proves \eqref{eq:joint-uniform-nondegeneracy}. Uniformity is obtained
at this last step, not by assuming uniform constants for measurable
phase regularity or for the product set.
\end{proof}

\begin{remark}[The role of the section]
The algebra in \eqref{eq:joint-zero-coboundary} uses only that the
section indicator is integer-valued and that its mean lies strictly
between zero and one. The two rotation products force an integral
constant phase without an arithmetic assumption on this mean. The
specific section geometry is used earlier to establish the bounded
collision covariance and its $L^2$ kernel characterization. Thus the
present proof neither evaluates a discontinuous section weight on
stable arcs nor replaces it by a smoother section.
\end{remark}

\begin{corollary}[Actual covariances and the physical-time projection]
\label{cor:actual-uniform-ellipticity}
For all sufficiently large $n$, uniformly in $R$,
\begin{equation}\label{eq:finite-n-joint-ellipticity}
 \frac1n\int U_{n,R}\otimes U_{n,R}\dd\nu_R^*
       \ge\frac{d_0}{2}I_4.
\end{equation}
The physical-time fluctuation covariance $\mathcal V_R$ in
Corollary~\ref{cor:physical-functional-gaussian} is also uniformly
positive definite on its three-dimensional space.
\end{corollary}
\begin{proof}
Equation~\eqref{eq:actual-covariance-rate} and
\eqref{eq:joint-uniform-nondegeneracy} give
\eqref{eq:finite-n-joint-ellipticity}. For the physical covariance recall
$\Gamma_R=c_*D_R$ and
$\mathcal V_R=\bar\tau_R^{-1}B_R\Gamma_RB_R^{\mathsf T}$, where
$B_R(x_1,x_2,x_3,x_4)=(x_1,x_2,x_3-x_4/\bar\tau_R)$.
For $z\in\R^3$, $|B_R^{\mathsf T}z|\ge|z|$. Hence
$z^{\mathsf T}\mathcal V_Rz\ge
 c_*d_0(\max_I\bar\tau_R)^{-1}|z|^2$. All constants are positive.
\end{proof}

\subsection{The Gaussian part of exact local inversion}
\begin{corollary}[Uniform Gaussian inversion and the remaining terms]
\label{cor:elliptic-central-inversion}
The centered Gaussian density associated with the actual return record is
\begin{equation}\label{eq:uniform-joint-gaussian-density}
 g_R(\xi)=\frac{1}{(2\pi)^2\sqrt{\det D_R}}
              \exp\!\left(-\tfrac12\xi^{\mathsf T}D_R^{-1}\xi\right).
\end{equation}
It is defined for every $R\in I$. There are $C,c>0$ such that
\begin{equation}\label{eq:elliptic-gaussian-tail}
 \sup_R\int_{|v|>n^{1/200}}
          e^{-v^{\mathsf T}D_Rv/2}\dd v
       \le C\exp(-c n^{1/100}).
\end{equation}
In Corollary~\ref{cor:raw-from-growing-arc}, uniform positive definiteness
is therefore a proved conclusion for this physical family, not an
additional hypothesis. Under its stated residual integrability,
complementary-transform estimate and local-edge conditions, the raw
mixed-density conclusion follows with the density
\eqref{eq:uniform-joint-gaussian-density}.
\end{corollary}
\begin{proof}
Positive definiteness gives the Gaussian inversion formula. In four
dimensions, integration in polar coordinates gives, for $A>0$,
\[
 \int_{|v|>A}e^{-d_0|v|^2/2}\dd v
   =2\pi^2\left(\frac{A^2}{d_0}+\frac{2}{d_0^2}\right)
                         e^{-d_0A^2/2}.
\]
Absorb the polynomial into a weaker exponential and put $A=n^{1/200}$.
The last assertion follows from the exact decomposition and the proved
central comparison in Corollary~\ref{cor:raw-from-growing-arc}.
\end{proof}

The two Fourier tails in this last argument concern different objects.
Equation~\eqref{eq:elliptic-gaussian-tail} estimates the limiting Gaussian.
It does not estimate the complementary transform of the physical raw
record or of its edge-subtracted residual. The proved central band is
still $|v|\le2n^{1/200}$, equivalently
$|\omega|\le2n^{-99/200}$. The annulus beyond it, the full critical and
singular edge extraction, and the weighted tails and relative event
comparison for exact physical conditioning remain the analytical tasks
in the raw mixed-density theorem. The covariance obstruction has been
removed without changing that theorem's record or target measure.
